\chapter{Muestra y protocolo de revisi\'on interna}
\label{cap:anexod}

\section*{D.1. Finalidad del anexo}

Este anexo documenta el instrumento preparado para revisar recomendaciones P1--P4
sobre casos reales del Registro 112 Castilla y Le\'on. Su funci\'on es habilitar el
bucle humano: los casos de riesgo y las discrepancias deben poder examinarse,
justificarse y registrarse de forma trazable.

La muestra interna no equivale a una validaci\'on externa completada. Se han
generado los casos y los campos de registro para que puedan revisarse de forma trazable,
pero esta parte no sustituye una validaci\'on formal con operadores del 112.

\section*{D.2. Muestra de revision}

La seleccion de casos combina dos criterios:

\begin{itemize}
    \item inclusion de los 59 falsos negativos P1 detectados en el test estratificado;
    \item inclusion de una muestra equilibrada adicional de 60 casos, con 15 casos por
    clase P1, P2, P3 y P4.
\end{itemize}

El total resultante es de 119 casos revisables. La muestra supera el m\'inimo previsto
y sirve tambi\'en como entrada de la evaluaci\'on automatizada de fidelidad. Esta
segunda evaluaci\'on se ejecuta mediante un juez determinista local sobre
explicaciones de Capa 3 en modo degradado.

\begin{table}[ht]
\centering
\caption{Composici\'on de la muestra preparada para revisi\'on}
\label{tab:anexod-muestra}
\small
\begin{tabular}{|l|r|}
\hline
\textbf{Tipo de caso} & \textbf{Numero de casos} \\ \hline
Falsos negativos P1 & 59 \\ \hline
Muestra equilibrada P1 & 15 \\ \hline
Muestra equilibrada P2 & 15 \\ \hline
Muestra equilibrada P3 & 15 \\ \hline
Muestra equilibrada P4 & 15 \\ \hline
\textbf{Total} & \textbf{119} \\ \hline
\end{tabular}
\end{table}

\section*{D.3. Plantilla de revision humana}

Cada fila del fichero de validacion representa un caso revisable. La plantilla incluye
identificacion, referencia academica, prediccion del sistema, evidencias mostradas y
campos reservados para la decision humana.

\begin{table}[ht]
\centering
\caption{Campos principales de la plantilla de revision}
\label{tab:anexod-plantilla}
\scriptsize
\begin{tabularx}{\textwidth}{|p{3.5cm}|X|}
\hline
\textbf{Campo} & \textbf{Funcion} \\ \hline
Identificador de revision & Codigo interno de caso, con formato REV-NNN \\ \hline
Tipo de muestra & Falso negativo P1 o muestra equilibrada \\ \hline
Identificador del incidente & Referencia trazable del Registro 112 CyL \\ \hline
Texto operativo & Texto minimizado del incidente, usado para revisar la decision \\ \hline
Etiqueta academica & Prioridad P1--P4 generada por supervision debil \\ \hline
Recomendacion del sistema & Prioridad P1--P4 recomendada por RuleFit-lite \\ \hline
Probabilidades & Distribucion normalizada por clase \\ \hline
Acuerdo local & Consistencia entre fuentes debiles \\ \hline
Reglas o se\~nales mostradas & Evidencias principales de la clase predicha \\ \hline
Contexto de explicacion & Prioridad, margen, evidencias y avisos usados por Capa 3 \\ \hline
Decision del revisor & Aceptar, modificar o rechazar la recomendacion \\ \hline
Prioridad final del revisor & Prioridad asignada por el revisor interno \\ \hline
Motivo de divergencia & Justificacion obligatoria si se modifica o rechaza \\ \hline
Auditoria especial & Marca casos que requieren revision prioritaria \\ \hline
\end{tabularx}
\end{table}

\section*{D.4. Casos reales seleccionados}

La Tabla~\ref{tab:anexod-casos} muestra una muestra compacta de casos incluidos en el
CSV. Los casos REV-061--REV-068 son falsos negativos P1 y, por tanto, tienen prioridad
de revision.

\begin{table}[ht]
\centering
\caption{Ejemplos de casos reales seleccionados para validacion}
\label{tab:anexod-casos}
\scriptsize
\begin{tabular}{|l|l|l|l|r|l|l|}
\hline
\textbf{Caso} & \textbf{Tipo} & \textbf{Incidente} & \textbf{Weak label} & \textbf{Sistema} & \textbf{Conf.} & \textbf{Categoria} \\ \hline
REV-001 & Equilibrado P1 & 1228180006005 & P1 & P1 & 0.712510 & trafico \\ \hline
REV-016 & Equilibrado P2 & 1264142203135 & P2 & P2 & 0.673363 & sanitario \\ \hline
REV-031 & Equilibrado P3 & 1284226557820 & P3 & P3 & 0.953758 & sanitario \\ \hline
REV-046 & Equilibrado P4 & 1284154490226 & P4 & P4 & 0.985624 & sanitario \\ \hline
REV-061 & FN P1 & 1244182009991 & P1 & P3 & 0.893565 & incendio \\ \hline
REV-062 & FN P1 & 1248678215637 & P1 & P2 & 0.552789 & incendio \\ \hline
REV-064 & FN P1 critico & 1255644701359 & P1 & P4 & 0.985624 & sanitario \\ \hline
REV-068 & FN P1 critico & 1279887542362 & P1 & P4 & 0.985624 & sanitario \\ \hline
REV-069 & FN P1 critico & 1284166596701 & P1 & P4 & 0.985624 & sanitario \\ \hline
REV-080 & FN P1 critico & 1284268940694 & P1 & P4 & 0.857180 & sanitario \\ \hline
REV-098 & FN P1 critico & 1284746951651 & P1 & P4 & 0.985624 & sanitario \\ \hline
REV-118 & FN P1 critico & 1285206245620 & P1 & P4 & 0.985624 & sanitario \\ \hline
\end{tabular}
\end{table}

\section*{D.5. Procedimiento de revision}

Cada autor deber\'a revisar de forma independiente los casos asignados, consultando el texto
operativo, la etiqueta debil, la recomendacion del sistema, las probabilidades, las
se\~nales activadas y, cuando este disponible, la explicacion normativa de Capa 3. La
decisi\'on se registrar\'a en tres valores posibles:

\begin{itemize}
    \item \textbf{Aceptar}: la recomendacion del sistema se considera adecuada.
    \item \textbf{Modificar}: el revisor asigna otra prioridad P1--P4 y justifica el
    cambio.
    \item \textbf{Rechazar}: el caso revela una recomendacion no util o evidencia
    insuficiente para sostenerla.
\end{itemize}

Las divergencias de dos o m\'as niveles se marcar\'an como auditor\'ia especial. En particular,
los casos P1 clasificados como P3 o P4 deber\'an analizarse antes de una futura validaci\'on externa,
porque representan infraestimaciones relevantes. En la version evaluada aparecen 7
casos P1--P4, que requieren especial atencion por representar infrapriorizaciones
criticas. Su analisis debe comprobar si el error procede de la guia de etiquetado, de
una extraccion incompleta de senales, de la ponderacion del modelo o de una falta de
evidencia explicita en el texto operativo.

\section*{D.6. Evaluaci\'on exploratoria con participantes}

Como complemento a la evaluacion tecnica, se realizo una evaluaci\'on breve del producto m\'inimo viable
(MVP) con nueve participantes. El acceso se facilito mediante un enlace de prueba basado en Google Colab y se
acompa\~no de casos representativos y un cuestionario estructurado sobre
comprensibilidad, utilidad percibida, confianza y claridad de la supervisi\'on humana.

La seleccion se distribuye entre tres perfiles. Para preservar la privacidad,
los participantes se identifican mediante c\'odigos A1--C3 y descripciones
funcionales, sin recoger nombres, unidades concretas ni datos personales. La
Tabla~\ref{tab:anexod-participantes} resume el dise\~no de la muestra.

\begin{table}[ht]
\centering
\caption{Perfiles anonimizados de la evaluaci\'on exploratoria del MVP}
\label{tab:anexod-participantes}
\scriptsize
\begin{tabularx}{\textwidth}{|p{0.9cm}|p{3.2cm}|X|p{3.9cm}|}
\hline
\textbf{Codigo} & \textbf{Grupo} & \textbf{Perfil funcional} & \textbf{Rol en la evaluacion} \\ \hline
A1 & Puesto de mando/CIS/emergencias & Perfil de transmisiones y sistemas de informaci\'on y comunicaciones (CIS), con experiencia en coordinaci\'on de puestos de mando. & Coherencia operativa, flujo de informaci\'on y trazabilidad. \\ \hline
A2 & Puesto de mando/CIS/emergencias & Perfil vinculado al montaje de puestos de mando en emergencias. & Utilidad en contexto de emergencia y claridad de prioridad cr\'itica. \\ \hline
A3 & Puesto de mando/CIS/emergencias & Perfil de inform\'atica y redes en puestos de mando. & Trazabilidad t\'ecnica, continuidad y comprensi\'on del sistema. \\ \hline
B1 & Tecnico IA & Tecnico de telecomunicaciones radio, usuario habitual de IA. & Claridad tecnica y utilidad en comunicaciones radio. \\ \hline
B2 & Tecnico IA & Tecnico de telecomunicaciones satelite, usuario habitual de IA. & Claridad en conectividad y sistemas distribuidos. \\ \hline
B3 & Tecnico IA & Ingeniero informatico y desarrollador, usuario habitual de IA. & Arquitectura, separacion de capas y reproducibilidad. \\ \hline
C1 & No especializado IA & Tecnico de montaje de puestos de mando, sin uso tecnico habitual de IA. & Comprension desde perfil tecnico no IA. \\ \hline
C2 & No especializado IA & Usuario no tecnico, sin uso habitual de IA. & Comprensibilidad general. \\ \hline
C3 & No especializado IA & Perfil de seguridad publica, usuario no tecnico de IA. & Percepcion desde seguridad publica y supervision humana. \\ \hline
\end{tabularx}
\end{table}

El cuestionario utiliza una escala Likert de 1 a 5 y preguntas abiertas. La
Tabla~\ref{tab:anexod-resultados-usuarios} resume los resultados agregados por grupo.

\begin{table}[ht]
\centering
\caption{Resultados agregados de la evaluaci\'on exploratoria}
\label{tab:anexod-resultados-usuarios}
\scriptsize
\begin{tabular}{|p{4.1cm}|r|r|r|r|}
\hline
\textbf{Dimension evaluada} & \textbf{Grupo A} & \textbf{Grupo B} & \textbf{Grupo C} & \textbf{Media global} \\ \hline
Prioridad recomendada razonable & 4.47 & 4.13 & 3.83 & 4.14 \\ \hline
Explicacion clara & 4.17 & 4.37 & 3.60 & 4.04 \\ \hline
Se\~nales/reglas aumentan confianza & 4.50 & 4.23 & 3.40 & 4.04 \\ \hline
Decision humana clara & 4.93 & 4.43 & 4.53 & 4.63 \\ \hline
\end{tabular}
\end{table}

Los comentarios abiertos muestran tres patrones. Los perfiles vinculados a puestos de
mando valoran la trazabilidad y la supervision humana; los usuarios tecnicos piden
mantener clara la separacion entre modelo, reglas y explicacion; y los usuarios no
especializados solicitan explicaciones mas directas de probabilidades y se\~nales. Esta
lectura orienta mejoras de interfaz y comunicacion, pero no debe interpretarse como
validacion operativa.

\section*{D.7. Limitaciones}

La muestra procede de etiquetas academicas generadas por supervision debil. Por tanto,
la revision interna puede detectar incoherencias del sistema o de la propia etiqueta,
pero no certifica validez operativa. La validacion externa con personal del 112 sigue
siendo una linea prioritaria de trabajo futuro.

La evaluaci\'on exploratoria con nueve participantes permite obtener indicios
descriptivos sobre utilidad percibida y comprension, pero no conclusiones
estadisticamente generalizables ni validacion operativa del sistema.
